\title{Group Sequence Policy Optimization}

\begin{abstract}
We present Group Sequence Policy Optimization (GSPO), a reinforcement learning algorithm developed to efficiently and reliably train large language models. GSPO departs from traditional methods that utilize importance ratios calculated at the token level; instead, it formulates importance ratios over the likelihood of entire sequences and applies sequence-level mechanisms for clipping, reward assignment, and optimization. Our findings indicate that GSPO not only surpasses the GRPO algorithm in terms of training speed and model quality but also brings enhanced stability to Mixture-of-Experts (MoE) reinforcement learning training. Furthermore, GSPO can streamline reinforcement learning system architectures. These advantages of GSPO have played a key role in the significant advancements observed in the recent Qwen3 models.
\end{abstract}

\section{Introduction}
\label{sec:intro}

Reinforcement learning (RL) has become a central framework for advancing the scalability of language models \citep{o1,dpsk-r1,qwq32b,qwen3}. Leveraging large-scale RL allows these models to acquire the ability to solve intricate tasks, including competitive mathematics and programming challenges, by engaging in extended, complex reasoning.

Achieving effective scaling of RL requires ensuring consistent and resilient training behaviors as computational resources increase. Unfortunately, modern RL techniques, such as GRPO \citep{grpo}, are plagued by substantial stability problems when applied to enormous language models, frequently causing drastic, irrecoverable breakdowns in the models \citep{qwen3, minimax-m1}. This lack of stability significantly impedes advancements in language model capabilities through further RL training.

This paper demonstrates that GRPO's instability originates from a critical error in the use and validity of importance sampling weights within its algorithm. This flaw generates substantial variance in training noise, which intensifies as response length grows and is exacerbated by the clipping strategy, leading ultimately to model degradation.

To overcome these inherent weaknesses, we introduce \textbf{Group Sequence Policy Optimization (GSPO)}, a novel reinforcement learning algorithm tailored for training large language models. GSPO's central contribution is a theoretically robust formulation of importance ratios grounded in sequence likelihood \citep{click}, thereby adhering more closely to the core principles of importance sampling. Moreover, GSPO calculates normalized rewards as the advantage values across multiple responses to a single query, facilitating proper alignment between reward assignment at the sequence level and the overall optimization process.

Our experiments clearly indicate that GSPO outperforms GRPO across key dimensions such as training stability, computational efficiency, and overall performance. Importantly, GSPO intrinsically addresses the instability issues typically encountered during reinforcement learning with large Mixture-of-Experts (MoE) architectures, thus removing the requirement for elaborate stabilization methods and offering opportunities to streamline RL systems. These advantages have played a central role in achieving substantial gains within the most recent Qwen3 models. Looking ahead, we see GSPO as a strong and extensible basis that can support the future growth of large-scale reinforcement learning with language models.

\section{Preliminaries}

\paragraph{Notation}
Throughout this work, we refer to an autoregressive language model with parameters $\theta$ as the policy $\pi_\theta$. The variable $x$ represents a query, while $\mathcal{D}$ is used for the set of queries. For any response $y$ given a query $x$, the probability assigned by the policy $\pi_\theta$ is expressed as $\pi_\theta (y | x) = \prod_{t=1}^{|y|} \pi_\theta (y_t | x, y_{<t})$, where $|y|$ is the total number of tokens in $y$. The pair $(x, y)$, consisting of a query and its corresponding response, can be evaluated by a verifier $r$, which assigns a reward $r(x, y) \in [ 0, 1 ]$.

\paragraph{Proximal Policy Optimization (PPO)} Using samples generated from the old policy $\pi_{\theta_\text{old}}$, PPO \citep{ppo} constrains the policy update within a proximal region of the old policy through the clipping mechanism. Specifically, PPO employs the following objective for policy optimization (we omit the KL regularization term hereinafter for brevity, as it is not the focus of this paper):

\begin{align}
\mathcal{J}_\text{PPO}(\theta) = \mathbb{E}_{ x \sim \mathcal{D},\, y \sim \pi_{\theta_\text{old}}( \cdot | x) }
\left[ \frac{1}{|y|} \sum_{t=1}^{|y|} 
\min \left( w_{t}(\theta) \widehat{A}_{t},  \, \mathrm{clip} \left( w_{t}(\theta), 1 - {\varepsilon}, 1 + {\varepsilon}\right) \widehat{A}_{t} \right)
\right],
\end{align}

Here, the importance ratio for a token $y_t$ is expressed as $
w_{t}(\theta) = \frac{ \pi_{\theta} (y_{t} | x, y_{<t}) }{ \pi_{\theta_\text{old}} (y_{t} | x,y_{<t})} $. The advantage $\widehat{A}_{t}$ associated with $y_t$ is derived from a separate value model, and the parameter $\varepsilon$ denotes the clipping threshold for the importance ratios.

A primary obstacle in applying PPO arises from its substantial dependence on the value model. Typically, the value model is constructed to be nearly as large as the policy model, resulting in significant demands for both memory and computation. Additionally, the success of the PPO algorithm is tightly coupled with the precision of the value estimate. Building a trustworthy value model is not straightforward; extending its reliability to handle longer sequences and increasingly complex tasks becomes particularly problematic.

\paragraph{Group Relative Policy Optimization (GRPO)} GRPO \citep{grpo} bypasses the need for the value model by computing the relative advantage of each response within a group of responses to the same query. Specifically, GRPO optimizes the following objective:

\begin{align}
\label{equ:grpo}
\mathcal{J}_\text{GRPO}(\theta) = \mathbb{E}_{ x \sim \mathcal{D},\, \{y_i\}_{i=1}^G \sim \pi_{\theta_\text{old}}( \cdot | x) }
\left[ \frac{1}{G} \sum_{i=1}^{G} \frac{1}{|y_i|} \sum_{t=1}^{|y_i|} 
\min \left( w_{i,t}(\theta) \widehat{A}_{i,t},  \, \mathrm{clip} \left( w_{i,t}(\theta), 1 - {\varepsilon}, 1 + {\varepsilon}\right) \widehat{A}_{i,t} \right)
\right],
\end{align}

where $G$ is the number of generated responses to each query $x$ (i.e., the group size), and the importance ratio $w_{i,t}(\theta)$ and advantage $\widehat{A}_{i,t}$ of token $y_{i,t}$ are:

\begin{align}
    w_{i,t}(\theta)=\frac{ \pi_{\theta} (y_{i,t} | x, y_{i,<t}) }{ \pi_{\theta_\text{old}} (y_{i,t} | x,y_{i,<t})},\quad\ 
    \widehat{A}_{i,t} = \widehat{A}_{i} = \frac{r(x, y_i) - \mathrm{mean} \left( \{ r(x, y_i) \}_{i=1}^G \right) }{ \mathrm{std} \left( \{ r(x, y_i) \}_{i=1}^G \right) },
\end{align}

respectively, where all the tokens in $y_i$ share the same advantage as $\widehat{A}_{i}$.

\section{Motivation}
\label{sec:motivation}

As models increase in size and complexity—including the adoption of sparsity structures such as those found in Mixture-of-Experts architectures—and as their outputs lengthen, achieving efficient hardware use during RL training often requires substantial rollout batch sizes. To enhance the efficiency with which data is leveraged, it is common to subdivide these extensive batches into smaller mini-batches for subsequent gradient calculations. However, this batching process creates an off-policy learning scenario: the generated responses $y$ are drawn from a prior policy $\pi_{\theta_\text{old}}$, rather than the policy $\pi_\theta$ that is currently undergoing optimization. This dynamic is a key reason for the integration of clipping strategies in PPO and GRPO, which restrict the influence of highly ``off-policy'' samples in the gradient update process.

While mechanisms like clipping aim to manage this off-policy discrepancy, we identify a more fundamental issue in GRPO: \textit{its objective is ill-posed}. This problem becomes particularly acute when training large models on long-response tasks, leading to catastrophic model collapse. The ill-posed nature of the GRPO objective stems from a misapplication of importance sampling weights. The principle of importance sampling is to estimate the expectation of a function $f$ under a target distribution $\pi_\text{tar}$ by re-weighting samples drawn from a behavior distribution $\pi_\text{beh}$:

\begin{align}
\mathbb{E}_{z \sim \pi_\text{tar}} \left[ f(z) \right] = \mathbb{E}_{z \sim \pi_\text{beh}} \left[ \frac{ \pi_\text{tar} (z) }{ \pi_\text{beh} (z)} \, f(z) \right].
\end{align}

This approach fundamentally depends on averaging across a large number of samples ($N \gg 1$) drawn from the behavioral distribution $\pi_\text{beh}$, thereby allowing the importance weighting $\frac{ \pi_\text{tar} (z) }{ \pi_\text{beh} (z)}$ to accurately compensate for the disparity between the target and behavior distributions.

By contrast, GRPO incorporates the importance weight $\frac{ \pi_{\theta} (y_{i,t} | x, y_{i,<t}) }{ \pi_{\theta_\text{old}} (y_{i,t} | x,y_{i,<t})}$ at each individual token position $t$. Because this weighting is computed from only one sampled token $y_{i,t}$ of the subsequent-token distribution $\pi_{\theta_\text{old}}( \cdot | x, y_{i,<t})$, it does not effectively achieve the goal of correcting for off-policy bias. Rather, it introduces substantial variance into the gradient estimates during training; this variance is amplified over longer generated sequences and further intensified by the application of clipping. Our empirical findings reveal that this process frequently results in catastrophic model degradation, from which recovery is rarely possible. If this collapse occurs, attempts at remediation—such as restoring an earlier checkpoint, carefully adjusting hyperparameters (including clipping limits), increasing sequence length, or altering RL sampling strategies—have consistently failed to recover the model.

These results indicate a deeper problem within the GRPO framework. The inadequacy of token-level importance weighting highlights a key guideline: \textbf{the granularity of the optimization target must be consistent with the granularity at which the reward is provided}. Since reward is given at the sequence level, token-wise off-policy correction is fundamentally misaligned. This insight encourages us to abandon token-level objectives in favor of importance weighting and optimization methods that operate at the \textit{sequence level}.

\section{Algorithm}

\subsection{GSPO: Group Sequence Policy Optimization}

Although the use of token-level importance weights $\frac{ \pi_{\theta} (y_{i,t} | x, y_{i,<t}) }{ \pi_{\theta_\text{old}} (y_{i,t} | x,y_{i,<t})}$ poses challenges within GRPO, our analysis indicates that, within language generation tasks, employing the \textit{sequence-level} importance weight $\frac{ \pi_{\theta} (y | x) }{ \pi_{\theta_\text{old}} (y | x)}$ is theoretically sound: it quantifies the extent to which a response $y$, generated according to $\pi_{\theta_\text{old}} (\cdot | x)$, diverges from the updated policy $\pi_{\theta} (\cdot | x)$. This interpretation is consistent with the sequence-level reward structure and provides an informative perspective for understanding the rationale behind the clipping mechanism.

Based on this straightforward observation, we propose the \textbf{Group Sequence Policy Optimization (GSPO)} algorithm. GSPO employs the following sequence-level optimization objective:

\begin{align}
\mathcal{J}_\text{GSPO} (\theta) =
\mathbb{E}_{ x \sim \mathcal{D},\, \{y_i\}_{i=1}^G \sim \pi_{\theta_\text{old}}( \cdot | x) }
\left[ 
\frac{1}{G} \sum_{i=1}^{G}
\min \left( s_{i}(\theta)  \widehat{A}_{i},  \, \mathrm{clip} \left( s_{i}(\theta), 1 - {\varepsilon}, 1 + {\varepsilon} \right) \widehat{A}_{i} \right) 
\right],
\label{equ:gspo}
\end{align}

where we adopt the group-based advantage estimation:

\begin{align}
\widehat{A}_{i} = \frac{r(x, y_i) - \mathrm{mean} \left( \{ r(x, y_i) \}_{i=1}^G \right) }{ \mathrm{std} \left( \{ r(x, y_i) \}_{i=1}^G \right) },
\end{align}

and define the importance ratio $s_{i}(\theta)$ based on sequence likelihood \citep{click}:

\begin{align}
s_{i}(\theta) = \left( \frac{ \pi_{\theta} (y_i | x) }{ \pi_{\theta_\text{old}} (y_i | x)} \right)^{\frac{1}{|y_i|}}
=
\exp \left( \frac{1}{|y_i|} \sum_{t=1}^{|y_i|} \log \frac{ \pi_{\theta} (y_{i,t} | x, y_{i,<t}) }{ \pi_{\theta_\text{old}} (y_{i,t} | x,y_{i,<t})} \right).
\end{align}

Consequently, GSPO performs clipping at the level of whole responses rather than at individual tokens, thereby removing samples that are excessively ``off-policy'' from gradient computation. This strategy aligns with both sequence-level reward assignment and the optimization process. Additionally, we implement length normalization for $s_{i}(\theta)$ to help minimize variance and to keep $s_{i}(\theta)$ within a consistent numerical interval. If omitted, the shift in probability for just a few tokens could cause significant volatility in the sequence-level importance ratios, and responses of varying lengths would necessitate different clipping thresholds. Furthermore, we observe that the clipping bounds used in GSPO are generally several orders of magnitude different from those employed in earlier methods such as GRPO, owing to the disparate ways importance ratios are defined.

\subsection{Gradient Analysis}
\label{subsec:gradient}

We can derive the gradient of the GSPO objective as follows (clipping is omitted for brevity):

\begin{align}
\nabla_{\theta} \mathcal{J}_\text{GSPO} (\theta)
=&\ 
\nabla_{\theta} \mathbb{E}_{ x \sim \mathcal{D},\, \{y_i\}_{i=1}^G \sim \pi_{\theta_\text{old}}( \cdot | x) }
\left[ 
\frac{1}{G} \sum_{i=1}^{G} s_{i}(\theta) \widehat{A}_{i}
\right] \\
= &\ 
\mathbb{E}_{ x \sim \mathcal{D},\, \{y_i\}_{i=1}^G \sim \pi_{\theta_\text{old}}( \cdot | x) }
\left[ 
\frac{1}{G} \sum_{i=1}^{G}
s_{i}(\theta) \widehat{A}_{i}
\cdot \nabla_{\theta} \log s_{i}(\theta)
\right] \\
=&\ 
\mathbb{E}_{ x \sim \mathcal{D},\, \{y_i\}_{i=1}^G \sim \pi_{\theta_\text{old}}( \cdot | x) }
\left[ 
\frac{1}{G} \sum_{i=1}^{G} 
\left( \frac{ \pi_{\theta} (y_i | x) }{ \pi_{\theta_\text{old}} (y_i | x)} \right)^{\frac{1}{|y_i|}} \widehat{A}_{i} 
\cdot \frac{1}{|y_i|} \sum_{t=1}^{|y_i|} \nabla_{\theta} \log \pi_{\theta} (y_{i,t} | x, y_{i,<t}) 
\right].
\label{equ:gspo_grad}
\end{align}

For comparison, the gradient of the GRPO objective is as follows (note that $\widehat{A}_{i,t} = \widehat{A}_{i}$):

\begin{align}
\nabla_{\theta} \mathcal{J}_\text{GRPO} (\theta) 
=&\ 
\nabla_{\theta} \mathbb{E}_{ x \sim \mathcal{D},\, \{y_i\}_{i=1}^G \sim \pi_{\theta_\text{old}}( \cdot | x) }
\left[ \frac{1}{G} \sum_{i=1}^{G} \frac{1}{|y_i|} \sum_{t=1}^{|y_i|} 
w_{i,t}(\theta) \widehat{A}_{i,t}
\right] \\
=&\
\mathbb{E}_{ x \sim \mathcal{D},\, \{y_i\}_{i=1}^G \sim \pi_{\theta_\text{old}}( \cdot | x) }
\left[ \frac{1}{G} \sum_{i=1}^{G} \widehat{A}_{i} 
\cdot \frac{1}{|y_i|} \sum_{t=1}^{|y_i|} 
\frac{ \pi_{\theta} (y_{i,t} | x, y_{i,<t}) }{ \pi_{\theta_\text{old}} (y_{i,t} | x,y_{i,<t})} 
\nabla_{\theta} \log \pi_{\theta} (y_{i,t} | x, y_{i,<t})  
\right].
\end{align}

Thus, the primary difference between GSPO and GRPO pertains to \textit{the approach they take toward weighting gradients of token log likelihoods}. GRPO assigns weights to tokens based on their individual ``importance weights'' $\frac{ \pi_{\theta} (y_{i,t} | x, y_{i,<t}) }{ \pi_{\theta_\text{old}} (y_{i,t} | x,y_{i,<t})}$. These varying weights—ranging from $(0, 1+\varepsilon]$ when $\widehat{A}_{i} >0$, or $[1-\varepsilon, +\infty)$ when $\widehat{A}_{i} < 0$—cannot be disregarded; their effects may compound over the course of training, potentially resulting in unintended and erratic outcomes. Conversely, GSPO applies uniform weighting to every token in a generated response, thereby removing the source of instability introduced by GRPO.

\subsection{GSPO-token: A Token-level Objective Variant}

In scenarios like multi-turn RL, we may desire a finer-grained advantage adjustment than the sequence level. To this end, we introduce a token-level objective variant of GSPO, namely \textbf{GSPO-token}, to allow token-wise advantage customization:

\begin{align}
\mathcal{J}_\text{GSPO-token}(\theta)
=
\mathbb{E}_{ x \sim \mathcal{D},\, \{y_i\}_{i=1}^G \sim \pi_{\theta_\text{old}}( \cdot | x) }
\left[ 
\frac{1}{G} \sum_{i=1}^{G} \frac{1}{|y_i|} \sum_{t=1}^{|y_i|} 
\min \left( s_{i,t}(\theta) \widehat{A}_{i,t},  \, \mathrm{clip} \left( s_{i,t}(\theta), 1 - {\varepsilon}, 1 + {\varepsilon} \right) \widehat{A}_{i,t} \right) 
\right],
\label{equ:gspo-token}
\end{align}

where

\begin{align}
s_{i,t}(\theta) 
= \mathrm{sg} \left[ s_{i}(\theta) \right]  \cdot \frac{ \pi_{\theta} (y_{i,t} | x, y_{i,<t}) }{ \mathrm{sg} \left[ \pi_{\theta} (y_{i,t} | x, y_{i,<t}) \right] },
\end{align}

and $\mathrm{sg}[\cdot]$ denotes only taking the numerical value but stopping the gradient, corresponding to the \texttt{detach} operation in PyTorch. The gradient of GSPO-token can be derived as:

\begin{align}
\nabla_{\theta} \mathcal{J}_\text{GSPO-token}(\theta)
=&\ 
\nabla_{\theta} \mathbb{E}_{ x \sim \mathcal{D},\, \{y_i\}_{i=1}^G \sim \pi_{\theta_\text{old}}( \cdot | x) }
\left[ 
\frac{1}{G} \sum_{i=1}^{G}
\frac{1}{|y_i|} \sum_{t=1}^{|y_i|} 
s_{i,t}(\theta) \widehat{A}_{i,t}
\right] \\
=&\ 
\mathbb{E}_{ x \sim \mathcal{D},\, \{y_i\}_{i=1}^G \sim \pi_{\theta_\text{old}}( \cdot | x) }
\left[ 
\frac{1}{G} \sum_{i=1}^{G} s_i (\theta) \cdot \frac{1}{|y_i|} \sum_{t=1}^{|y_i|} 
\widehat{A}_{i,t} \frac{ \nabla_{\theta} \pi_{\theta} (y_{i,t} | x, y_{i,<t}) }{ \pi_{\theta} (y_{i,t} | x, y_{i,<t}) }
\right] \\
=&\ 
\mathbb{E}_{ x \sim \mathcal{D},\, \{y_i\}_{i=1}^G \sim \pi_{\theta_\text{old}}( \cdot | x) }
\left[ 
\frac{1}{G} \sum_{i=1}^{G}  \left( \frac{ \pi_{\theta} (y_i | x) }{ \pi_{\theta_\text{old}} (y_i | x)} \right)^{\frac{1}{|y_i|}}
\cdot \frac{1}{|y_i|} \sum_{t=1}^{|y_i|}
\widehat{A}_{i,t} \nabla_{\theta} \log \pi_{\theta} (y_{i,t} | x, y_{i,<t})  
\right].
\label{equ:gspo-token_grad}
\end{align}

It is important to recognize that the expression $\frac{ \pi_{\theta} (y_{i,t} | x, y_{i,<t}) }{ \mathrm{sg} \left[ \pi_{\theta} (y_{i,t} | x, y_{i,<t}) \right] }$ evaluates to 1, which implies that $s_{i,t}(\theta)$ is numerically equivalent to $s_{i}(\theta)$. By examining Equation~\eqref{equ:gspo} alongside Equation~\eqref{equ:gspo-token}, as well as Equation~\eqref{equ:gspo_grad} and Equation~\eqref{equ:gspo-token_grad}, it becomes evident that GSPO-token and GSPO yield identical numerical results for the optimization objective, clipping criteria, and theoretical gradient whenever all tokens in a response $y_i$ are assigned a uniform advantage (specifically, when $\widehat{A}_{i,t} = \widehat{A}_{i}$ for every $t$). Nevertheless, GSPO-token introduces enhanced versatility, as it allows the assignment of distinct advantage values to each token individually.

\section{Experiments and Discussion}

\subsection{Empirical Results}

We conduct experiments utilizing a cold-start model that is fine-tuned from Qwen3-30B-A3B-Base, presenting both training reward trajectories and model performance metrics across several benchmarks: AIME'24 (reporting average Pass@1 over 32 samples), LiveCodeBench (202410-202502, average Pass@1 computed from 8 samples), and CodeForces (Elo Rating). Throughout RL training, each batch of rollout data is divided into four mini-batches for the purpose of applying gradient updates. For GSPO, we assign the left and right clipping bounds in Equation~\eqref{equ:gspo} to values of 3e-4 and 4e-4, correspondingly. Our comparisons involve GRPO as a baseline, where the left and right clipping values in Equation~\eqref{equ:grpo} are set to 0.2 and 0.27—hyperparameters carefully chosen to ensure fairness. It should be emphasized that GRPO requires the adoption of the Routing Replay training approach to achieve stable MoE RL convergence (further explored in \S~\ref{subsec:routing_replay}), whereas \textit{GSPO eliminates the dependency on this technique}.

As illustrated in Figure~\ref{fig:results}, the GSPO training process maintains stability from start to finish. We find that \textbf{GSPO supports ongoing enhancements in performance by increasing computational resources during training, periodically refreshing the query set, and lengthening the generated sequences}. Additionally, GSPO demonstrates higher efficiency compared to GRPO, reaching superior accuracy and benchmark results given equivalent computational budgets and query usage. In conclusion, GSPO has been successfully implemented in RL training of state-of-the-art Qwen3 models, convincingly validating GSPO’s effectiveness in leveraging RL scaling for large language models.

\subsection{Curious Observation on Clipping Fractions}

An important difference between GSPO and GRPO lies in how responses are clipped: GSPO clips at the response level, whereas GRPO clips individual tokens. As illustrated in Figure~\ref{fig:clipping}, the proportion of clipped tokens is nearly two orders of magnitude higher for GSPO than for GRPO, and this gap persists even when the clipping ranges are adjusted. Surprisingly, although GSPO clips many more tokens, resulting in fewer tokens available for training or gradient calculations, it outperforms GRPO in terms of training efficiency. This unexpected result — where more aggressive clipping actually enhances efficiency — suggests that the gradient estimates in GRPO, which are computed at the token level, suffer from significant noise and inefficiency in leveraging samples. GSPO’s sequence-level strategy, by contrast, delivers a clearer and more potent signal for learning.

\subsection{Benefit of GSPO for MoE Training}
\label{subsec:routing_replay}

\paragraph{Background}
In contrast to RL training with dense architectures, MoE models present distinct challenges concerning training stability due to their sparse activation patterns. Notably, our investigations indicate that employing the GRPO algorithm can result in pronounced \textit{expert-activation volatility} within MoE models, impeding proper convergence during RL optimization. Specifically, we observe that the set of experts engaged for identical responses may shift dramatically following one or more gradient update steps. For instance, using the 48-layer Qwen3-30B-A3B-Base configuration, about 10\% of the experts activated for any given rollout example differ between the updated policy $\pi_{\theta}$ and the previous policy $\pi_{\theta_\text{old}}$ after each RL gradient update. This effect intensifies in larger and deeper MoE models, causing sharp and erratic variations in the token-level importance ratios $w_{i,t}(\theta) = \frac{ \pi_{\theta} (y_{i,t} | x, y_{i,<t}) }{ \pi_{\theta_\text{old}} (y_{i,t} | x,y_{i,<t})}$, as explored in \S~\ref{sec:motivation} and~\ref{subsec:gradient}. These large oscillations undermine the validity of the importance weights, ultimately disrupting stable RL convergence.

\paragraph{Our Previous Approach} In earlier work, we addressed this problem by adopting the \textbf{Routing Replay} training method. The approach involves storing the expert routing choices from $\pi_{\theta_\text{old}}$ and subsequently using these identical routes when evaluating $\pi_{\theta}$ for computing importance ratios: $w_{i,t}(\theta) = \frac{ \pi_{\theta} (y_{i,t} | x, y_{i,<t}) }{ \pi_{\theta_\text{old}} (y_{i,t} | x,y_{i,<t})}$. As a result, both $\pi_{\theta}$ and $\pi_{\theta_\text{old}}$ utilize the same expert configuration for each token $y_{i,t}$, ensuring that the calculated importance ratios remain stable at the token level. This shared activation also facilitates consistent optimization of the network's routing decisions during training updates. Figure~\ref{fig:routing_replay} illustrates how Routing Replay is pivotal for maintaining standard convergence behavior within GRPO training for MoE architectures.

\paragraph{Benefit of GSPO} While Routing Replay facilitates proper convergence in GRPO training for MoE models, its reliance on fixed routing strategies leads to greater memory usage and additional communication requirements. Moreover, it can artificially restrict the capacity available in the MoE architecture. By contrast, as illustrated in Figure~\ref{fig:results}, GSPO operates without dependence on Routing Replay. It allows standard computation of importance ratios $s_i(\theta)$, achieving reliable convergence and stable optimization. The principal observation is that GSPO attends exclusively to the likelihood of the entire sequence, $\pi_{\theta} (y_i | x)$, as opposed to being influenced by token-level probabilities $\pi_{\theta} (y_{i,t} | x, y_{i,<t})$. Given that the language modeling function of the MoE model remains intact, the sequence likelihood exhibits minimal variability. Therefore, GSPO offers a definitive solution to the problem of fluctuating expert activation seen in MoE models, eliminating the necessity for intricate approaches such as Routing Replay. This streamlines the training, reinforces stability, and enables utilization of the model’s complete capacity without imposed limitations.

\subsection{Benefit of GSPO for RL Infrastructure}

Due to existing differences in numerical precision between engines used for training (such as Megatron) and those used for inference (like SGLang and vLLM), it is standard practice to employ the training engine to re-evaluate the likelihoods of generated responses under the previous policy $\pi_{\theta_\text{old}}$. However, because GSPO optimizes based on the sequence-level likelihoods instead of token-level likelihoods, it inherently exhibits greater robustness to such precision variations. As a result, GSPO enables the direct utilization of likelihood values provided by the inference engine for optimization purposes, thus removing the necessity of recalculating them with the training engine. This feature is particularly advantageous in contexts such as partial rollout, multi-turn RL, or disaggregated training and inference pipelines.

\section{Conclusion}

This paper introduces Group Sequence Policy Optimization (GSPO), a novel reinforcement learning approach tailored for the training of large language models. Rooted in the principle of importance sampling, GSPO utilizes sequence likelihood to define importance ratios, applying sequence-level mechanisms for clipping, reward assignment, and optimization. Empirical results show that GSPO offers significantly enhanced stability, training efficiency, and overall performance when compared to GRPO, and is especially effective for large-scale reinforcement learning involving MoE models. This methodology has played a critical role in achieving the remarkable enhancements present in the latest Qwen3 models. By positioning GSPO as a scalable foundational algorithm, this work aims to further advance RL research and anticipates substantial progress in the development of intelligent systems.

\end{document}